\section{Introduction}
% Stakes
A mobile agent that cooks from a recipe app and then fills a grocery list reads ``vanilla extract'' on screen 1 and must type it on screen 9, after the screen that held the text has left the display. We call a step whose required text appeared on an earlier screen, and only there, a \emph{delayed information need}; the screen that held the text is the \emph{source screen}, and the number of screens since is the \emph{gap}. Such steps can decide whether a long cross-app task completes, and the benchmarks used to develop mobile agents contain few of them. AITZ~\citep{aitz} has none in 18,643 steps, consistent with tasks that average 7.4 steps; GUI-Odyssey~\citep{guiodyssey}, at 15.4 steps per task, has at least one in 31\% of sampled trajectories (\S\ref{sec:needs}).

% Problem gap
Two prompted designs recover such text, and both degrade as the history grows. Keeping the full OCR text of every earlier screen in context, which we call \emph{full history}, localizes the right screen in 68\% of needs that arise by step 15 and in 46\% of later ones. A gist-and-lookup \emph{screen index}~\citep{readagent}, one short line per screen from which a reader picks a screen to re-read, localizes 41\% overall against full history's 54.8\%. Five ways of training or re-prompting its writer leave it there, and an oracle over sampled lines reaches full-history level and no higher (\S\ref{sec:index}). Prompting is the shared constraint. The one trained look-back agent we know of, PAL-UI~\citep{palui}, needs 8.6K curated step-level instructions, so training either design on a new corpus means paying for labels of which earlier screen to re-read.

% Key abstraction
Nobody has to. A string check defines a delayed information need, and that check also answers the training question. For a verified need at step $t$, every earlier screen either contains the typed text in its OCR or does not, so the verifier that mines the need labels all $t$ candidates at once. We call training a look-back policy on these labels \emph{verifier-supervised look-back}. The labels cost one substring test per screen, and no person inspects any of them.

% Design intuition
We train two instantiations on the verifier's labels. They differ in how the history reaches the model and, as a consequence, in how many decisions each label set yields. The \emph{per-screen scorer} is a 7B model with a low-rank adapter. It reads the task, the current screen, and one earlier screen at a time, and outputs the probability that this earlier screen holds the needed text; localization is the argmax over screens. The \emph{full-history reader} is the same base model fine-tuned on the same labels to emit a screen number from the whole history in one call. A second call, the \emph{read stage}, then types the text, either from the chosen screen alone (\emph{single-screen read}) or from the full history with the chosen screen repeated and marked (\emph{read with history}); \emph{end-to-end} accuracy scores the typed text.

% Results preview
Supervision, not architecture, carries the result. On 188 verified cases in five folds grouped by trajectory, training lifts the full-history reader from 54.8\% to 65.4\% ($p=0.019$) and the per-screen scorer from 50.0\% to 72.3\% ($p<0.001$). The two trained variants differ by 6.9 points of localization and 2.1 end-to-end, and a paired test separates neither ($p=0.086$ and $p=0.45$). The scorer's gain concentrates at gaps of 8--12 screens, where it localizes 78.2\% against 56.4\% ($p=0.003$). We mined a fourth sample after fixing every design decision and evaluated on it once. There the scorer localizes 79.6\% against 76.3\% for the trained full-history reader ($p=0.63$) and 59.1\% for prompted full history ($p=0.003$); read with history raises end-to-end accuracy from 23.7\% to 49.5\% for the scorer and 51.6\% for the trained full-history reader ($p<0.001$ for both). Both trained variants fall below prompted full history at gaps of 13 or more screens.

% Contributions
We contribute: (1) \textbf{Problem and labels.} A definition of delayed information needs, an automatic verifier that both mines and labels them, and their prevalence, with two filters any study of them requires (\S\ref{sec:needs}). (2) \textbf{Method.} Verifier-supervised look-back and two instantiations that share one objective (\S\ref{sec:method}). (3) \textbf{Finding.} The supervision is the lever: it lifts both architectures, and the per-screen decomposition adds a gain confined to mid-range gaps (\S\ref{sec:main}, \S\ref{sec:gap}). (4) \textbf{Validation and limits.} An external sample confirms the margin over prompting, ablations show the design is insensitive to its prompt and loss settings, and gaps of 13 or more screens remain open (\S\ref{sec:ext}, \S\ref{sec:limits}).
